\section{Detailed Methods and Protocol History}
\label{app:detailed-methods}

\begin{figure}[t]
\centering
\includegraphics[width=\linewidth]{figures/study_design.pdf}
\caption{Controlled experiment design and collected response counts. For each main
question, we test five models in four student-work and feeling conditions,
with two requests per combination. Main training and test source families are separate.
This diagram covers the first stage (5,320 responses); the subsequent expression
extension adds 5,120 responses on the same 32 test sources.
RQ1 and the first-stage RQ2 test use the 32 test problems under the standard prompt; RQ3 reuses 16 of them
under additional prompt conditions. Repetitions and prompt variants do not add
independent sources. We compare four nested predictors and a predictor based on training-response
style, with the robustness checks described below.}
\label{fig:study-design}
\end{figure}

We define \emph{educational personality} as recurring teaching behavior and
patterns in how a model responds to student cues. Let $Y_{mqsar}^{(k)}$ denote whether model configuration $m$ expresses
behavior $k$ on source problem $q$, student state $s$, instruction condition $a$, and
independent request $r$. A default profile summarizes behavior under the
standard tutoring prompt over a specified distribution of educational inputs. A
conditional profile additionally represents how that behavior varies with
student state. We test both descriptions through predictions on separate source
problems. This definition does not require a fixed cross-context model
ranking, and it does not treat any identifiable linguistic signature as a
validated personality construct.

The unit $m$ is a model configuration, including its requested and returned
model names, provider, observation period, and generation settings. The controlled experiments request MiniMax-M3, MiniMax-M2.7, GLM-5.3,
DeepSeek-V4-Pro, and Doubao-Seed-2.0-Lite. In the measurement pilot, the request
alias GLM-5.2 returned GLM-5.3 throughout; those responses are identified by the
returned name and are not pooled with historical GLM-5.2 results. A provider may change model weights without changing the model name. Our
results therefore apply to the versions observed during the experiment.




\subsection{Separating Historical Analysis from New Tests}

The archive audit distinguishes raw records from successful deduplicated
responses, and distinguishes tutors from solvers and evaluators. Repeated
questions and conversations can occur under different benchmark names and
prompt variants. Consequently, benchmark names and local item identifiers are
not enough to identify independent samples. We keep related problems in
the same group when testing prediction on new inputs. Analyses informed by previous archive results remain
exploratory and are distinguished from predictions fixed before new outputs
were collected.

The deduplicated archive contains 899,025 records when different input
versions are kept separate. Reviewing current benchmark implementations
identifies 68,799 records in nine tutoring settings; most other records involve
evaluation or answering test questions. Classifying a task as tutoring does not show that its prompts were neutral
or comparable across models. Appendix~\ref{app:archive-scope} reports
the complete role inventory and a hash-verified inspection of the 122 tutor
response files, including missing information about how responses were generated.

We group related questions into \emph{source families}, treating potentially
related questions as one group. The controlled experiment uses 32 training
questions in 30 source families and 32 test questions in 32 separate families. Each
partition has eight questions in each of mathematics, science, language, and
logic/data reasoning. Two training question pairs share an averaging or unit-
conversion family and remain grouped in internal validation. We exclude the four pilot templates and versions that only substitute numbers.
The test evaluates new problems within these four subjects. It does not test
new subjects or establish that models have never encountered the concepts.

The questions, correct references, erroneous student attempts, and correct
unfinished work were authored by the research agent. Two model reviewers
checked all 64 sets of question materials before they were finalized,
yielding 128 complete reviews.
All five yes/no content checks passed in every review. Twenty-five reviews
also included comments, mostly about the deliberately incorrect or unfinished
student work. We retain these comments and the research agent's decisions
about them. This is content quality control without new
human annotation, not a claim of validation against human annotations.

\subsection{Student Work and Feeling Conditions}

Each main question has two student-work states, an erroneous attempt and correct
but unfinished work, crossed with two explicit affect statements, calm and
frustrated. These are four controlled input conditions. The work-state contrast
bundles correctness and progress; it does not isolate a pure progress or latent
student-ability effect. The affect statement is an input manipulation, not an
observation of a real learner's emotion.

All prospective main conditions provide the tutor with the same correct
instructor reference for that question, specifying that the student has not
seen it. This equalizes the availability of the solution but does not fully
control comprehension, attention, or instruction-following ability. Results
under this controlled reference condition are reported separately from the
historical archive.

The training stage uses the standard tutoring prompt on all 32 questions. The
test stage first repeats this design on its 32 new questions. Within
each test domain, four questions are selected by a fixed hash for four
additional prompt conditions: two neutral paraphrases, an instruction to prioritize
student reasoning without revealing the final answer, and an instruction to
explain the solution and state the answer. Each selected question is tested under every instruction and student state. The neutral prompt conditions change only the role wording; the two
instruction conditions explicitly change the requested teaching behavior.

Every model--input combination receives two requests at temperature 0.7 with an
8,192-token output budget. The fixed design specifies 1,280 training responses,
1,280 test responses to the standard prompt, and 2,560 additional intervention
responses. It also includes 160 responses to eight paired test problems with complete
or missing information. Four training problems are repeated as reference
checks across stages, adding 40 responses. These checks and prompt variants
do not increase the number of independent test problems. Request manifests, input hashes,
parameters, timestamps, and failure histories are recorded. Requests are
interleaved within shared provider concurrency limits rather than run as
separate model blocks.

\subsection{Observable Events and Measurement Development}

We selected three primary actions using the measurement-development pilot.
\emph{Answer revelation} records an explicit final answer, including an incorrect answer or a correct
statement that the available information is insufficient.
It therefore measures disclosure, not correctness. \emph{Reasoning elicitation}
records an invitation for the student to complete a reasoning step, including
direct requests and indirect suggestions; rhetorical questions immediately answered
by the tutor do not qualify. \emph{Affect acknowledgement} records explicit
recognition of feeling or difficulty, excluding generic politeness and praise.
These are behavior measures, not three assumed independent personality factors.

Five secondary events cover worked explanation, learner choice, epistemic
qualification, requests for missing facts, and unsupported ability attribution.
All are retained regardless of their results. The pilot showed why a high
aggregate agreement rate is insufficient: requesting missing facts had only one
majority-coded positive example among 160 responses, while unsupported ability
attribution had weak positive agreement. Neither is promoted to a primary
outcome based on subsequent test results.

MiniMax-M3, DeepSeek-V4-Pro, and GLM-5.3 each annotate every tutor response.
Annotators receive the student conversation, reference answer, and response
with its original line numbers. They do not receive the generating model name,
system-role assignment, repeat number, or partition label. Natural output style
can still reveal identity. Positive codes must reference existing nonempty
source lines; negative codes have no references. Valid line references make the judgment traceable but do not ensure it is
correct. We use majority labels when all three annotation outputs are valid.
We also analyze each annotator separately and exclude each in turn, assigning
0.5 to ties between two annotators.

The initial exact-quotation protocol failed on naturally formatted answers,
which motivated line references during measurement development. Original
failures remain archived. The finalized pilot has 480 valid codes for 160
unchanged tutor responses after two structure/completion failures were retried
under identical inputs. Formal structure failures are likewise recorded and
retried without selecting by event label or annotator agreement.
A training coding interruption was recovered using saved valid outputs and
the existing two-pass repair allowance. This yielded 3,839 of 3,840 valid codes.
One GLM coding request repeatedly exhausted its 16,384-token completion budget
without visible output. Before fixing predictions or test generation, a
disclosed amendment allowed this single missing code a 32,768-token ceiling
and at most two attempts, keeping its model, input, rubric, and temperature
fixed. It completed on the first attempt using 3,511 completion tokens; this
does not identify the cause of the preceding failures. The original and
expanded-budget groups were validated separately before their disjoint codes
were combined. All 1,280 training answers consequently have three valid codes.
The interrupted run, failed attempts, and amendment receipt are retained; the
single expanded-budget code is an explicit deviation from the original
uniform-budget protocol.

Test coding subsequently exhausted its two repair passes with 12,114
of 12,120 codes valid. All six remaining GLM requests produced empty visible
outputs after both 16,384-token attempts in each repair pass. Before computing
primary test results, a separate amendment permitted these six requests
a 32,768-token ceiling and at most two attempts each, with unchanged inputs,
model, rubric, temperature, and parser. It selects every remaining missing code,
without selecting by event label or agreement; all prior valid codes are retained.
We validate the two budget groups separately and combine them only after
checking that no annotation is duplicated or omitted.
Two affected answers were generated under the standard prompt for prediction tests, one under a neutral
wording condition, and three in the tests with missing information. The
supplemental sensitivity enumerates all binary assignments for the two missing
GLM labels for responses to the standard prompt while keeping the other annotations and predictions unchanged, and
bounds the affected descriptive contrasts. This is an explicitly amended
measurement protocol; it does not reset the original frozen repair allowance.
Five requests completed under this amendment. The remaining opportunity-probe
request returned Hypertext Transfer Protocol (HTTP) status 502 at approximately
300 seconds on both attempts,
matching the local Gateway's upstream read timeout. A separate transport
amendment retained those five completed codes, increased only the idle
Gateway's read timeout from 300 to 660 seconds, and permitted at most two
additional attempts for that one request under the same 32,768-token budget.
The client timeout remained 720 seconds. This extra transport allowance is
also retained in the protocol history rather than counted as part of either
earlier repair allowance.


\subsection{Fixing Predictions Before Collecting Test Responses}
\label{app:prediction-method}

The predictors are statistical logistic-regression models. The LLMs generate
tutor answers and code their behaviors; they do not produce the primary
behavior predictions. For each event, training responses from the five models
are pooled, with model identity included as an input feature. The
predictors use structured subject and student-state indicators rather than
reading the full problem text. Fitting a predictor does not train or fine-tune
any evaluated LLM. The same five models are evaluated on the separate
test questions.

For each behavior, four nested logistic predictors model the probability of its
occurrence: a subject-by-student-state baseline; that baseline plus model
identity; an additional model-by-subject block; and an additional
model-by-student-state block. New feature blocks are orthogonalized against
existing blocks using training data and standardized under source-equal
weights. This prevents duplicate main-effect encodings from changing their
ridge penalty when interactions are introduced. The intercept is unpenalized.
The penalty is selected from $\{0.1,1,10,100\}$ by source-grouped training-only
validation. Constant training outcomes use the same smoothed constant across
all four predictors.

The fitted models, tuning choices, input hashes, and test probabilities
are saved before any test tutor response is generated. The generation
runner checks this prediction lock. No test label, answer length, or
same-item cross-model score enters a primary predictor. A secondary competing
predictor uses only training-source profiles of log word count and marked-line
fraction, recomputing those profiles within training validation. Such style
competition is descriptive: teaching policy may itself influence length and
formatting, so it is not a causal adjustment for style.

For an action with predicted probability $p$ and observed binary label
$y\in\{0,1\}$, Brier loss is $(p-y)^2$. A positive gain is the baseline's
average loss minus the augmented predictor's average loss, so it indicates
improved probability prediction.

For each of the three primary events, the two primary comparisons are the Brier
improvement from a default model profile over the situation baseline and from a
student-state profile over the model-by-subject baseline. We average losses
within source and then equally across the 32 test sources. The six
one-sided paired source-level tests use Holm correction. We report absolute
Brier losses, source-bootstrap intervals, and all comparison signs. These
intervals condition on the fitted training profiles and a source-sampling
approximation; they do not describe uncertainty over all models or all
educational contexts.

\subsection{Repetition and Prompt Plasticity}

Repetition analyses report positive, negative, and total event agreement and
profile correspondence, retaining source clustering. Constant behavior can
produce perfect repeat agreement and is reported as such. On the 16 fully
crossed intervention sources, we estimate each model's event-rate difference
between an alternative role and the standard tutoring prompt. We report average shifts
alongside model and student-state heterogeneity.

For neutral wording, the fixed equivalence reference is $\pm0.10$ in event
probability. We retain the nominal Bonferroni paired-$t$ intervals across the
30 model--event--neutral-prompt comparisons. A nonsignificant difference does not
establish equivalence, and constant source differences receive a degeneracy
flag. Further synthetic diagnostics during training show undercoverage in a
sparse boundary case even with nonzero source variance. Before test
generation, we therefore additionally require a bounded-source interval to
support any population-equivalence statement. Scaling Hoeffding's bound to
source differences in $[-1,1]$ and allocating error across both tails and all
30 comparisons gives half-width
$\sqrt{2\log(2\cdot30/0.05)/n}$ \citep{hoeffding1963}. At $n=16$ this is
approximately $0.941$, too wide to support the chosen tolerance. This
conservative sensitivity still assumes independent source observations and
does not make the authored bank representative. We can describe observed
neutral shifts and nominal intervals, while withholding population equivalence.
These disclosed interpretation additions preserve the frozen numerical output
and the six primary prediction tests. The
explicit ASK and EXPLAIN conditions measure changes in the requested actions, without assessing teaching quality. Finally, answer revelation and reasoning elicitation are analyzed
jointly as four possible combinations. Their evidence lines do not identify
which action occurred first, and we do not infer action order from them.

\subsection{Sensitivity Analyses and Their Timing}

During training coding, after measurement development but before test
responses exist, we additionally fix eight alternative sets of annotations: the
three-annotator mean, each individual annotator, each two-annotator mean, and the mean after
excluding annotators from the generating model's family. For each label set, predictors are refitted
using its training labels and locked before test generation. They are
evaluated against the corresponding test labels; two-annotator ties remain
0.5. This checks dependence on coding choices but cannot remove biases shared
by all annotators or supply an external ground truth.

Five further analyses each exclude one generating model from both training
and testing. We refit and tune the predictors using the remaining training data. This checks dependence on a single configuration;
it changes the target model mixture and does not test prediction of an unseen
model. Neither sensitivity family replaces the six primary comparisons.

We also resample the 30 training source families 200 times, preserving the
multiplicity of repeated source draws. We retain the selected regularization
relative to total loss weight, without repeating hyperparameter selection, and
recompute training-style profiles in each resample. The resulting gain quantiles
describe training-source sensitivity conditional on the test material
and selected regularization. They are separate from the primary test-
source intervals. These extensions were specified before test outputs,
not preregistered before the entire research project began.

\subsection{Mathematics and Content-Error Sensitivity}

The protocol also requires a mathematics subset and a check of dependence on
obvious content errors. Its implementation is fixed during training coding,
before test generation. We evaluate the existing predictions on all eight
mathematics sources, covering 320 responses to the standard prompt, without refitting.
Separately, MiniMax-M3 and DeepSeek-V4-Pro review all 1,280 responses to the standard prompt
for checkable content errors using the student input, reference, and original
answer lines. They distinguish a clear error, no clear error identified, and
uncertainty. Withholding a solution, asking a question, or correctly quoting
and correcting an error is not itself an error. Positive error judgments
require a category and source-line evidence.

Twelve control examples written by research agents accompany the responses, giving 2,584
secondary review requests. Each reviewer must match at least five of six
controls in each of the clear-error and no-clear-error classes to meet the
limited diagnostic rule. If either reviewer fails, screened results remain
measurement-limited diagnostics. These controls do not establish correctness
on natural responses, even when the rule is met.

We retain the full primary results and report three additional screens:
excluding unanimously identified errors; excluding errors identified by either
reviewer; and excluding either errors or uncertainty. If a response is flagged, we remove responses from all five models for the
same problem, student state, and repetition, keeping the comparisons matched. We report retained coverage and source-equal losses
with descriptive source intervals, without new hypothesis tests. These screens
condition on generated responses and change the target distribution; they are
not causal controls for ability. The content review covers the responses to the standard prompt,
not all prompt-intervention responses.

The first content pass returned all 2,584 records, but its initial validator
reconstructed a zero temperature as a JavaScript Object Notation (JSON)
integer \texttt{0}, whereas the
request runner had serialized float \texttt{0.0}. A separately recorded
correction reconstructs the exact payload from the verified recorded numeric
value. All prompt, response and corrected payload hashes match; no sampling
parameter, input, verdict rule, or original file changes. The unchanged parser
accepts 2,577 initial reviews. The first of the original two allowed repair
passes resolves all seven structural failures: two invalid concise explanations
and five invalid evidence-line references. All seven complete on their first
HTTP attempts; the 2,577 valid original reviews remain unchanged. The correction,
original failed diagnostic, and repair history are retained. Both reviewers'
twelve control judgments are valid and match the expected verdicts.

\subsection{Implementation and Computational Reporting}
\label{app:implementation}
The saved analysis environment records Python 3.13.11, NumPy 2.5.2,
SciPy 1.18.0, pandas 3.0.5, scikit-learn 1.9.0, statsmodels 0.14.6,
Matplotlib 3.11.1, and seaborn 0.13.2. Statistical analyses run locally;
tutor generation and behavioral coding use hosted model APIs. No tutor
weights are trained for this study. Parameter counts and provider-side graphics processing unit (GPU)
hours of the served configurations were not independently verified.

Generation and coding counts, completion limits, retries, and missing-code
amendments are reported by stage above and in
Appendices~\ref{app:cue-transfer} and~\ref{app:archive-validation}.
These include 10,440 controlled tutor replies and 3,584 separately archived
replies, with pilot and measurement-development work reported separately.
The million-record archive predates the controlled experiments; its total
compute cost cannot be reconstructed from the retained provenance. We do not
convert request counts into unobserved training costs, energy, or carbon estimates.
